\documentclass{beamer}

\usepackage[utf8]{inputenc}
\input{../helper}

\title{Introdução à Computação \\ tipos de dados em C, entrada e saída}
\author{Alexandre Rademaker}
\date{}

\begin{document}

\begin{frame}
 \maketitle
\end{frame}

\begin{frame}[fragile,allowframebreaks]{variáveis, literais e tipos}

  \textbf{Variáveis} são nomes para localizações na memória que
  poderão conter valores e um valor é sempre de um determinado
  \textbf{tipo}.

  O tipo só precisa ser especificado na declaração da variável. São
  \textbf{palavras reservadas}.

  Literais são notações para representar um valor.\vspace{.5cm}

  \begin{lstlisting}[style=normalc]
    int a = 12, b;
    double pi = 3.14;
    float b = 1.3E2;

    char test;
    test = 'a';
  \end{lstlisting}

  \framebreak
  
  Em C, cada \href{https://en.wikipedia.org/wiki/C_data_types}{tipo de
    dado} é armazenado em uma certa quantidade de bytes.

  \begin{description}
  \item[bool] 1 byte
  \item[char] 1 byte
  \item[short] 2 bytes
  \item[int] 4 bytes (32 bits) que podem representar até $2^{32}$
    valores diferentes.
  \item[long] 8 bytes    
  \item[float] 4 bytes
  \item[double] 8 bytes    
  \end{description}

  entre outros \ldots

  \framebreak

  Os identicadores (nomes) não pode ser uma \textbf{palavra reservada}
  (\texttt{int}, \texttt{true}, \texttt{if}, \texttt{else} etc).

  São formados por uma combinação de letras, dígitos e '\_'
  (underline), começando sempre com uma letra ou '\_'.
  
  Letras maiúsculas e minúsculas são diferentes. para assegurar a
  portabilidade use no máximo 31 caracteres. escolha nomes
  \textbf{significativos} para clareza do código.

  \framebreak

  Alguns erros comuns:

  \begin{lstlisting}[style=normalc]
    int var1, 2var, _var3;
    int int;
    int double;
  \end{lstlisting}

  Variáveis podem ser declaradas em qualquer lugar de um programa C,
  mas devem aparecer antes de serem usadas no programa. Variáveis tem
  um escopo de vida, veremos mais à frente.

  \framebreak

  \begin{columns}
    \begin{column}{.5\textwidth}
      \begin{lstlisting}[style=normalc]
        char c = 'a';
        int b = 2630692;
      \end{lstlisting}

      \vspace{2cm}
      Variáveis tem nome, valor e endereço (acessível via
      {\color{red}\texttt{\&c}})
      
    \end{column}
    \begin{column}{.5\textwidth}
      \includegraphics[width=.8\textwidth]{vars-memory.png}
    \end{column}
  \end{columns}
  
\end{frame}


\begin{frame}[fragile]{operador de atribuição}

  \begin{lstlisting}[style=normalc]
    s = 10 + n;
  \end{lstlisting}

  \begin{itemize}
  \item O {\color{red} =} é atribuição, não teste de igualdade. O
    {\color{red} +} é a adição. Ambos binários.

  \item A expressão à direita é avaliação e seu resultado atribuído a
    variável {\color{red} s}.

  \item Sempre usar espaços entre operadores e operandos, ajuda a
    leitura.
  \end{itemize}
\end{frame}

\begin{frame}[fragile]{operadores aritméticos}

  \begin{center}
    \begin{tabular}{cc}
      \hline
      operação & operador\\
      \hline
      precedência & ( )\\
      multiplicação & *\\
      soma & +\\
      subtração & -\\
      divisão & /\\
      resto / módulo & \%\\
      \hline
    \end{tabular}
  \end{center}
  
  \vfill
  
  \begin{lstlisting}[style=normalc]
    int a = (12 + 4) * 4 % 2;
  \end{lstlisting}
  
\end{frame}


\begin{frame}{entrada e saída}

  Um programa precisa receber dados para processar (entrada) e
  produzir uma saída.

  Podemos passar dados por parâmetros na linha de comando, ler de um
  arquivo ou solicitar interativamente ao usuário.

  Para produzir dados, podemos imprimir na tela ou gravar dados em um
  arquivo.

  Entre outras formas!
  
\end{frame}


\begin{frame}[fragile,allowframebreaks]{a função \texttt{printf}}

  A função \ilcode{printf} é parte da biblioteca \ilcode{<stdio.h>}.

  \texttt{printf({\color{red} <format-string>}, {\color{green} e1}, {\color{green} e2}, ...)}

  onde $e1, e2,\ldots$ são expressão pode ser: um nome de variável, um
  literal ou expressões envolvendo nomes e literais.\vspace{.6cm}

  \begin{lstlisting}[style=normalc]
    printf("hello world"); // certo

    printf(12); // erro
    int v = 12; 
    print(v); // erro

    printf("%d", v); // certo
  \end{lstlisting}

  \framebreak

  Alguns possíveis 'especificadores' para o comando printf:
  
  \begin{itemize}
  \item \texttt{\color{red} \%d} int (número inteiro)
  \item \texttt{\color{red} \%ld} long long (número inteiro)
  \item \texttt{\color{red} \%f} float (ponto flutuante)
  \item \texttt{\color{red} \%lf} double (ponto flutuante)
  \item \texttt{\color{red} \%c} char (caractere)
  \item \texttt{\color{red} \%s} string (cadeia de caracteres)  
  \end{itemize}

  \framebreak

  \begin{lstlisting}[style=normalc]
    int a = 10;
    int b = 20;
    printf("Soma de %d e %d = %d\n", a, b, a + b);
  \end{lstlisting}

  Irá produzir

  \texttt{Soma de 10 e 20 = 30}

  \framebreak

  \begin{lstlisting}[style=normalc]
    #include <stdio.h>

    int main() {
      printf("%-3s %8s\n"  , "Var", "Val");
      printf("%-3s %8.2f\n", "x"  , 10.222);
      printf("%-3s %8.2f\n", "y"  , 20.33);
      printf("%-3s %8.2f\n", "z"  , 30E2);
    }
  \end{lstlisting}

  \begin{lstlisting}[style=code]
    Var      Val
    x      10.22
    y      20.33
    z    3000.00
  \end{lstlisting}

  \framebreak

  ainda temos os caracteres especias

  \begin{itemize}
  \item \verb|\n| quebra de linha, passa para a linha debaixo
  \item \verb|\t| tabulação horizontal
  \item \verb|\"| aspas duplas
  \item \verb|\'| aspas simples ou apóstrofo
  \item \verb|\\| barra invertida
  \item \verb|\a| beep 
  \end{itemize}

\end{frame}  


\begin{frame}[fragile,allowframebreaks]{a função \texttt{scanf}}

  A função \ilcode{scanf} também é parte da biblioteca \ilcode{<stdio.h>}.
  
  Utilizada para ler da entrada padrão (terminal). O scanf tem algumas
  (grandes) diferenças em relação ao printf.
  
  \begin{itemize}
  \item A função printf imprime texto e o \textbf{valor} de variáveis
  \item A função scanf \textbf{altera o conteúdo} das variáveis
  \item Alterar conteúdo equivale a \textbf{modificar o que está na
      memória}. Por esta razão, sempre passamos um \textbf{endereço de
      memória} para a função scanf.
  \end{itemize}

  \framebreak

  \texttt{scanf({\color{red} <format-string>}, {\color{green} addr1}, {\color{green} addr2}, ...)}

  onde \texttt{addr1}, \texttt{addr2},\ldots são endereços de
  memória.

  \vspace{.6cm}

  \begin{lstlisting}[style=normalc]
    int x;
    scanf("%d", &x);

    scanf(x); // erro
    scanf("%d", x); // erro
  \end{lstlisting}
    
  \framebreak

  \begin{lstlisting}[style=normalc]
    int n;
    float m;
    char c;
    scanf("%d ",  &n);
    scanf("%c ",  &c);
    scanf("/%f/", &m);
    printf("n:[%d] c:[%c] m:[%.3f] \n", n, c, m);
  \end{lstlisting}

  \begin{lstlisting}[style=code]
    > clang test.c && ./a.out
    123   a    /4.5/
    n:[123] c:[a] m:[4.500]
  \end{lstlisting}

  Veja \ilcode{man 3 scanf}

\end{frame}


\begin{frame}[fragile,allowframebreaks]{exemplos}

  \lstinputlisting[style=tinyc,caption=src/calculator1.c]{src/calculator1.c}

  Por agora, apenas lembre-se que ao passar variáveis para o
  \ilcode{scanf}, deve passar com o \ilcode{\&} na frente.

  \framebreak

  Vamos fazer uma média de 3 valores? O que iremos ver na saída?

  \lstinputlisting[style=normalc,caption=src/scores0.c]{src/scores0.c}

  \framebreak

  \includegraphics[width=.9\textwidth]{scores-warn.png}

  \framebreak

  Soma de 3 {\color{red} int}eiros divididos por um {\color{red}
    int}eiro é um {\color{red} int}eiro e o resto descartado.

  `printf` está esperando um \ilcode{int} (\ilcode{\%d}) em seu
  argumento \ilcode{format string} (string de formatação).

  Mudando \ilcode{\%d} para \ilcode{\%f} \ldots A função `printf`
  precisa saber como ler e escrever os bytes!

\end{frame}


\end{document}

%%% Local Variables:
%%% mode: latex
%%% TeX-master: t
%%% End:
